%! AP Statistics — Unit 4, Lesson 4.12 — Homework
\documentclass[10pt]{article}
\usepackage{apstats-article}
\usepackage{apstats-boxes}
\usepackage{grading}
\gradingassignment{apstats-u04-l12-hw}
\gradingcourse{advanced_statistics}

\begin{document}

\pageheader{Unit 4, Lesson 4.12}{Homework}
\namedateperiod

\vspace{0.1in}

\begin{practicebox}
\small
\textbf{Directions:} For each problem, show all formula and calculator setup.
Write complete-sentence interpretations wherever asked.

\vspace{10pt}
\textbf{1.} A standard die is rolled until a 1 appears.
Let $X =$ the number of rolls needed.
\begin{enumerate}[label=(\alph*), leftmargin=2em, itemsep=8pt]
  \item Check BITS and state the distribution.\\
        \gradedhide{\writeline\gradedright{\raisebox{-13pt}{\gradedblank[id=hw1a, type=rubric, points=3, display={B, I, T, S met (1 or not; independent rolls; until first 1; $p = 1/6$): $X \sim$ Geom$(1/6)$.}]{4.8in}{\gradedlinesheight{3}-5pt}}}\\[1pt]\writeline\\[1pt]\writeline}
  \item Compute $P(X = 2)$ and $P(X = 6)$ using the formula. Show setup.\\
        \writeline\\[1pt]\writeline\gradedhide{\hrulefill}\gradedright{\raisebox{-2pt}{$P(X{=}2) =$~\gradedblank[id=hw1b1, type=short, points=1, answer=0.1389, check=numeric, tolerance=0.002]{0.75in}{0.2in}\quad $P(X{=}6) =$~\gradedblank[id=hw1b2, type=short, points=1, answer=0.0670, check=numeric, tolerance=0.002]{0.75in}{0.2in}}}
  \item Find $\mu_X$ and interpret in context.\\
        \writeline\gradedhide{\hrulefill}\gradedright{\raisebox{-2pt}{\gradedblank[id=hw1c1, type=short, points=1, answer=6, check=numeric, tolerance=0, display={\textbf{6}}]{0.9in}{0.2in}}}\\[1pt]\gradedhide{\writeline\gradedright{\raisebox{-1pt}{\gradedblank[id=hw1c2, type=rubric, points=2, display={Average of 6 rolls to get the first 1.}]{2.2in}{0.22in}}}}
  \item Find $P(X > 6) = (1 - p)^6$. Interpret in one sentence.\\
        \writeline\gradedhide{\hrulefill}\gradedright{\raisebox{-2pt}{\gradedblank[id=hw1d1, type=short, points=1, answer=0.3349, check=numeric, tolerance=0.002]{0.9in}{0.2in}}}\\[1pt]\gradedhide{\writeline\gradedright{\raisebox{-1pt}{\gradedblank[id=hw1d2, type=rubric, points=2, display={33.5\%: more than 6 rolls.}]{2.2in}{0.22in}}}}
\end{enumerate}

\vspace{8pt}
\textbf{2.} A factory produces bolts. Each bolt independently has a 4\% chance
of being too large. Bolts are tested one at a time until the first oversized bolt
is found.
\begin{enumerate}[label=(\alph*), leftmargin=2em, itemsep=8pt]
  \item State the distribution of $X$ (the number of bolts tested).\\
        \makebox[5cm][l]{\hspace{0.7cm}\smash{\raisebox{-4pt}{\gradedblank[id=hw2a, type=short, points=1, answer=Geom(0.04), check=text]{3cm}{0.2in}}}}
  \item Compute $P(X = 10)$. Write the formula and answer.\\
        \writeline\\[1pt]\writeline\gradedhide{\hrulefill}\gradedright{\raisebox{-2pt}{\gradedblank[id=hw2b, type=short, points=2, answer=0.0277, check=numeric, tolerance=0.001]{0.9in}{0.2in}}}
  \item Find $\mu_X$ and interpret.\\
        \writeline\gradedhide{\hrulefill}\gradedright{\raisebox{-2pt}{\gradedblank[id=hw2c1, type=short, points=1, answer=25, check=numeric, tolerance=0]{0.9in}{0.2in}}}\\[1pt]\gradedhide{\writeline\gradedright{\raisebox{-1pt}{\gradedblank[id=hw2c2, type=rubric, points=2, display={On average, 25 bolts are tested to find the first oversized one.}]{5in}{0.22in}}}}
  \item Compute $P(X \le 5)$ using \texttt{geometcdf}.\\
        \writeline\gradedhide{\hrulefill}\gradedright{\raisebox{-2pt}{\gradedblank[id=hw2d, type=short, points=1, answer=0.1846, check=numeric, tolerance=0.002]{0.9in}{0.2in}}}
\end{enumerate}

\vspace{8pt}
\textbf{3.} (Binomial vs.\ Geometric) A researcher samples voters.
Each voter independently supports a candidate with probability 0.55.
\begin{enumerate}[label=(\alph*), leftmargin=2em, itemsep=8pt]
  \item The researcher interviews voters until finding the first supporter.
        Let $W =$ the number interviewed. State the distribution of $W$.\\
        \makebox[5cm][l]{\hspace{0.7cm}\smash{\raisebox{-4pt}{\gradedblank[id=hw3a1, type=short, points=1, answer=Geom(0.55), check=text]{3cm}{0.2in}}}} \quad $\mu_W = $ \makebox[3cm][l]{\hspace{1.6cm}\smash{\raisebox{-4pt}{\gradedblank[id=hw3a2, type=short, points=1, answer=1.82, check=numeric, tolerance=0.01]{1.4cm}{0.2in}}}}
  \item The researcher interviews exactly 10 voters.
        Let $V =$ the number of supporters. State the distribution of $V$.\\
        \makebox[5cm][l]{\hspace{0.7cm}\smash{\raisebox{-4pt}{\gradedblank[id=hw3b1, type=short, points=1, answer={B(10, 0.55)}, check=text]{3cm}{0.2in}}}} \quad $\mu_V = $ \makebox[3cm][l]{\hspace{1.6cm}\smash{\raisebox{-4pt}{\gradedblank[id=hw3b2, type=short, points=1, answer=5.5, check=numeric, tolerance=0]{1.4cm}{0.2in}}}}
  \item Explain the key difference between $W$ and $V$ in one or two sentences.\\
        \writeline\\[1pt]\gradedhide{\writeline\gradedright{\raisebox{-1pt}{\gradedblank[id=hw3c, type=rubric, points=2, display={$W$: trials until 1st success (random count); $V$: successes in fixed 10.}]{4.6in}{0.22in}}}}
\end{enumerate}

\vspace{8pt}
\textbf{4.} (AP-Style Free Response) \textit{A certain allele appears in 20\% of
a large population's individuals. A geneticist examines individuals one at a time
(independently) until finding the first one carrying the allele.
Let $X =$ the number of individuals examined.}
\begin{enumerate}[label=(\alph*), leftmargin=2em, itemsep=8pt]
  \item Define the distribution of $X$. Check BITS briefly.\\
        \gradedhide{\writeline\gradedright{\raisebox{-13pt}{\gradedblank[id=hw4a, type=rubric, points=3, display={$X \sim$ Geom$(0.20)$: carrier or not; independent individuals (large population); until the first carrier; same $p = 0.20$.}]{6.1in}{\gradedlinesheight{3}-5pt}}}\\[1pt]\writeline\\[1pt]\writeline}
  \item Compute $P(X = 4)$. Show formula setup.\\
        \writeline\\[1pt]\writeline\gradedhide{\hrulefill}\gradedright{\raisebox{-2pt}{\gradedblank[id=hw4b, type=short, points=2, answer=0.1024, check=numeric, tolerance=0.001]{0.9in}{0.2in}}}
  \item Find $\mu_X$ and interpret in context.\\
        \writeline\gradedhide{\hrulefill}\gradedright{\raisebox{-2pt}{\gradedblank[id=hw4c1, type=short, points=1, answer=5, check=numeric, tolerance=0]{0.9in}{0.2in}}}\\[1pt]\gradedhide{\writeline\gradedright{\raisebox{-1pt}{\gradedblank[id=hw4c2, type=rubric, points=2, display={On average, 5 individuals examined to find the first carrier.}]{4.2in}{0.22in}}}}
  \item Find $P(X \le 3)$. Write the calculator expression and the answer.\\
        \writeline\gradedhide{\hrulefill}\gradedright{\raisebox{-2pt}{\gradedblank[id=hw4d, type=short, points=1, answer=0.488, check=numeric, tolerance=0.002]{0.9in}{0.2in}}}
  \item Find $P(X > 5)$. Interpret in context.\\
        \writeline\gradedhide{\hrulefill}\gradedright{\raisebox{-2pt}{\gradedblank[id=hw4e1, type=short, points=1, answer=0.3277, check=numeric, tolerance=0.002]{0.9in}{0.2in}}}\\[1pt]\gradedhide{\writeline\gradedright{\raisebox{-1pt}{\gradedblank[id=hw4e2, type=rubric, points=2, display={32.8\%: more than 5 examined.}]{2.3in}{0.22in}}}}
\end{enumerate}

\vspace{8pt}
\textbf{5. Extension:} The geometric distribution is \textbf{memoryless}.
Verify numerically that $P(X > 5 \mid X > 2) = P(X > 3)$ when $p = 0.20$.
Show both calculations and explain what ``memoryless'' means in the context
of the geneticist problem.

\vspace{4pt}
\writeline\\[1pt]\writeline\\[1pt]\writeline

\end{practicebox}

\end{document}
